\documentclass[12pt]{article}

% -------------------- preamble --------------------
\usepackage{amsthm}
\newtheorem*{exercise}{Exercise}
\newtheorem*{solution}{Solution}
\usepackage{amsmath, amssymb}
\usepackage{physics}
\usepackage{geometry}
\usepackage{hyperref}
\usepackage{nicematrix}
\usepackage{empheq}
\usepackage{xcolor}
\usepackage{tikz}
\usetikzlibrary{positioning, angles, quotes}
\usetikzlibrary{shapes.geometric}

\newcommand*\widefbox[1]{\fbox{\hspace{2em}#1\hspace{2em}}}

\geometry{margin=1in}

\title{Lecture Notes: Tensor Calculus\\
\large Based on Youtube series by eigenchris}
\author{}
\date{}

% -------------------- document --------------------
\begin{document}
\maketitle
\tableofcontents
\newpage

% -------------------- input sections --------------------
\section{Multi-variable Calculus Review}

For this lecture notes on Tensor Calculus, there are some concepts and notes that we need to quickly review on multi-variable calculus, which
need to be clear and understood before proceeding.\\
Here we'll give a brief review of the following topics:

\begin{itemize}
    \item Single-variable Derivatives and chain rule;
    \item Single-variable Integrals and differentials (change of variable);
    \item Multi-variable Partial derivatives and chain rule;
    \item Gradient of a function;
    \item Directional Derivatives
    \item Multi-variable Integrals and differentials;
    \item Arch length, velocity vector tangent to a curve.
\end{itemize}
\subsection{Single-Variable Derivatives and Chain Rule}

The derivative of a function $f(x)$ with respect to $x$ measures the instantaneous rate of change:
\[
\dv{f}{x} = \lim_{h \to 0} \frac{f(x+h) - f(x)}{h}
\]

\textbf{Basic derivative rules:}
\begin{align*}
\dv{x}(x^n) &= nx^{n-1} \\
\dv{x}(e^x) &= e^x \\
\dv{x}(\sin x) &= \cos x \\
\dv{x}(\cos x) &= -\sin x \\
\dv{x}(\ln x) &= \frac{1}{x}
\end{align*}

\textbf{Sum and product rules:}
\begin{align*}
\dv{x}(f + g) &= \dv{f}{x} + \dv{g}{x} \\
\dv{x}(fg) &= f\dv{g}{x} + g\dv{f}{x}
\end{align*}

\textbf{Chain rule:} For a composition of functions $f(g(x))$:
\[
\dv{x}f(g(x)) = \dv{f}{g} \cdot \dv{g}{x}
\]

More explicitly, if $y = f(u)$ and $u = g(x)$, then:
\[
\dv{y}{x} = \dv{y}{u} \cdot \dv{u}{x}
\]

\textbf{Example:} Let $y = \sin(x^2)$. Setting $u = x^2$, we have $y = \sin(u)$:
\[
\dv{y}{x} = \cos(u) \cdot 2x = 2x\cos(x^2)
\]

\subsection{Single-Variable Integrals and Change of Variable}

The definite integral of $f(x)$ from $a$ to $b$ represents the signed area under the curve:
\[
\int_a^b f(x) \, dx
\]

\textbf{Fundamental Theorem of Calculus:} If $F(x)$ is an antiderivative of $f(x)$ (i.e., $\dv{F}{x} = f$), then:
\[
\int_a^b f(x) \, dx = F(b) - F(a)
\]

\textbf{Basic integrals:}
\begin{align*}
\int x^n \, dx &= \frac{x^{n+1}}{n+1} + C \quad (n \neq -1) \\
\int \frac{1}{x} \, dx &= \ln|x| + C \\
\int e^x \, dx &= e^x + C \\
\int \sin x \, dx &= -\cos x + C \\
\int \cos x \, dx &= \sin x + C
\end{align*}

\textbf{Change of Variable (Substitution):} To integrate $f(g(x))g'(x)$, substitute $u = g(x)$, so that $du = g'(x) \, dx$:
\[
\int f(g(x)) g'(x) \, dx = \int f(u) \, du
\]

For definite integrals, the limits change accordingly: if $u = g(x)$, then $x = a \to u = g(a)$ and $x = b \to u = g(b)$.

\textbf{Example:} Evaluate $\int x \cos(x^2) \, dx$. Let $u = x^2$, then $du = 2x \, dx$, so $x \, dx = \frac{1}{2} du$:
\[
\int x \cos(x^2) \, dx = \int \cos(u) \cdot \frac{1}{2} \, du = \frac{1}{2}\sin(u) + C = \frac{1}{2}\sin(x^2) + C
\]

\subsection{Multi-Variable Partial Derivatives and Chain Rule}

For a function $f(x, y)$ of multiple variables, the \textbf{partial derivative} with respect to one variable treats all other variables as constants.

\textbf{Notation:}
\[
\pdv{f}{x}, \quad \pdv{f}{y}, \quad f_x, \quad f_y
\]

\textbf{Definition:}
\[
\pdv{f}{x} = \lim_{h \to 0} \frac{f(x+h, y) - f(x, y)}{h}
\]

\textbf{Example:} For $f(x, y) = x^2y + \sin(xy)$:
\begin{align*}
\pdv{f}{x} &= 2xy + y\cos(xy) \\
\pdv{f}{y} &= x^2 + x\cos(xy)
\end{align*}

\textbf{Multivariable Chain Rule:} Consider $z = f(x, y)$ where $x = x(t)$ and $y = y(t)$ are functions of $t$. Then:
\[
\dv{z}{t} = \pdv{f}{x}\dv{x}{t} + \pdv{f}{y}\dv{y}{t}
\]

More generally, if $z = f(x, y)$ and $x = x(u, v)$, $y = y(u, v)$, then:
\[
\pdv{z}{u} = \pdv{f}{x}\pdv{x}{u} + \pdv{f}{y}\pdv{y}{u}
\]

\textbf{Example:} Let $z = x^2 + y^2$ with $x = r\cos\theta$ and $y = r\sin\theta$. Then:
\begin{align*}
\pdv{z}{r} &= \pdv{z}{x}\pdv{x}{r} + \pdv{z}{y}\pdv{y}{r} \\
&= (2x)(\cos\theta) + (2y)(\sin\theta) \\
&= 2r\cos^2\theta + 2r\sin^2\theta = 2r
\end{align*}

\textbf{Einstein Summation Convention:} Using indexed notation where $x^i$ represents the $i$-th coordinate, the chain rule can be written compactly as:
\[
\pdv{f}{u^j} = \pdv{f}{x^i}\pdv{x^i}{u^j}
\]
where summation over the repeated index $i$ is implied (no explicit $\sum$ symbol needed). This convention will be used extensively in tensor calculus.

\subsection{Gradient of a Function}

The \textbf{gradient} of a scalar function $f(x, y, z)$ is a vector field that points in the direction of greatest rate of increase of the function.

\textbf{Definition:} For a function $f: \mathbb{R}^n \to \mathbb{R}$, the gradient is:
\[
\grad f = \nabla f = \begin{pmatrix} \pdv{f}{x_1} \\ \pdv{f}{x_2} \\ \vdots \\ \pdv{f}{x_n} \end{pmatrix}
\]

In three dimensions:
\[
\grad f = \pdv{f}{x}\hat{i} + \pdv{f}{y}\hat{j} + \pdv{f}{z}\hat{k}
\]

\textbf{Geometric interpretation:}
\begin{itemize}
    \item $\grad f$ points in the direction of steepest ascent of $f$
    \item $|\grad f|$ gives the rate of change in that direction
    \item $\grad f$ is perpendicular to level curves (2D) or level surfaces (3D) of $f$
\end{itemize}

\textbf{Example:} For $f(x, y) = x^2 + y^2$:
\[
\grad f = 2x\hat{i} + 2y\hat{j} = 2(x\hat{i} + y\hat{j})
\]
At point $(1, 1)$: $\grad f = 2\hat{i} + 2\hat{j}$, pointing radially outward.

\textbf{Applications:}
\begin{itemize}
    \item \textit{Optimization}: Finding maxima/minima (critical points where $\grad f = 0$)
    \item \textit{Physics}: Conservative forces $\vec{F} = -\grad U$ (gravitational, electric fields)
    \item \textit{Machine learning}: Gradient descent for minimizing loss functions
    \item \textit{Fluid dynamics}: Pressure gradients drive fluid flow
\end{itemize}

\subsection{Directional Derivatives}

While partial derivatives measure the rate of change along coordinate axes, the \textbf{directional derivative} measures how fast a function changes in an arbitrary direction.

\textbf{Definition:} The directional derivative of $f$ at point $\vec{a}$ in the direction of unit vector $\vec{u}$ is:
\[
D_{\vec{u}}f(\vec{a}) = \lim_{h \to 0} \frac{f(\vec{a} + h\vec{u}) - f(\vec{a})}{h}
\]

This gives the instantaneous rate of change of $f$ as we move from $\vec{a}$ in direction $\vec{u}$.

\textbf{Key requirement:} $\vec{u}$ must be a \textit{unit vector} (i.e., $|\vec{u}| = 1$) for the directional derivative to represent the true rate of change per unit distance.

\textbf{Formula using gradient:} The directional derivative can be computed using the dot product:
\[
D_{\vec{u}}f = \grad f \cdot \vec{u} = |\grad f||\vec{u}|\cos\theta = |\grad f|\cos\theta
\]
where $\theta$ is the angle between $\grad f$ and $\vec{u}$.

\textbf{Important properties:}
\begin{itemize}
    \item \textbf{Maximum rate:} $D_{\vec{u}}f$ is maximized when $\vec{u}$ points in the direction of $\grad f$ (i.e., $\theta = 0$):
    \[
    \max_{\vec{u}} D_{\vec{u}}f = |\grad f|
    \]
    \item \textbf{Minimum rate:} $D_{\vec{u}}f$ is minimized when $\vec{u}$ points opposite to $\grad f$ (i.e., $\theta = \pi$):
    \[
    \min_{\vec{u}} D_{\vec{u}}f = -|\grad f|
    \]
    \item \textbf{Zero rate:} $D_{\vec{u}}f = 0$ when $\vec{u}$ is perpendicular to $\grad f$ (i.e., $\theta = \pi/2$), meaning along level curves/surfaces
\end{itemize}

\textbf{Example 1:} Find the directional derivative of $f(x, y) = x^2 + xy$ at $(1, 2)$ in the direction toward $(4, 6)$.

First, compute the gradient:
\[
\grad f = (2x + y)\hat{i} + x\hat{j} \quad \Rightarrow \quad \grad f(1,2) = 4\hat{i} + \hat{j}
\]

The direction vector from $(1, 2)$ to $(4, 6)$ is $\vec{v} = 3\hat{i} + 4\hat{j}$. Normalize it:
\[
\vec{u} = \frac{\vec{v}}{|\vec{v}|} = \frac{3\hat{i} + 4\hat{j}}{\sqrt{9 + 16}} = \frac{3}{5}\hat{i} + \frac{4}{5}\hat{j}
\]

The directional derivative is:
\[
D_{\vec{u}}f = \grad f \cdot \vec{u} = (4\hat{i} + \hat{j}) \cdot \left(\frac{3}{5}\hat{i} + \frac{4}{5}\hat{j}\right) = \frac{12}{5} + \frac{4}{5} = \frac{16}{5}
\]

\textbf{Example 2:} For $f(x, y) = x^2 + y^2$ at $(3, 4)$, find the direction of maximum increase and the maximum rate.

\[
\grad f = 2x\hat{i} + 2y\hat{j} \quad \Rightarrow \quad \grad f(3,4) = 6\hat{i} + 8\hat{j}
\]

The direction of maximum increase is along $\grad f$. The unit vector is:
\[
\vec{u}_{\text{max}} = \frac{6\hat{i} + 8\hat{j}}{\sqrt{36 + 64}} = \frac{6\hat{i} + 8\hat{j}}{10} = \frac{3}{5}\hat{i} + \frac{4}{5}\hat{j}
\]

The maximum rate of increase is:
\[
|\grad f(3,4)| = \sqrt{36 + 64} = 10
\]

\subsection{Multi-Variable Integrals and Differentials}

\textbf{Double and triple integrals} extend the concept of integration to functions of multiple variables, computing volumes under surfaces or higher-dimensional analogues.

\textbf{Double integral:} For a function $f(x, y)$ over region $R$:
\[
\iint_R f(x, y) \, dA = \int_a^b \int_{g_1(x)}^{g_2(x)} f(x, y) \, dy \, dx
\]
where $dA = dx \, dy$ is the area element.

\textbf{Triple integral:} For a function $f(x, y, z)$ over region $V$:
\[
\iiint_V f(x, y, z) \, dV
\]
where $dV = dx \, dy \, dz$ is the volume element.

\textbf{Change of Variables (Jacobian):} When transforming coordinates $(x, y) \to (u, v)$, the area element transforms according to:
\[
dx \, dy = |J| \, du \, dv
\]
where $J$ is the \textbf{Jacobian determinant}:
\[
J = \pdv{(x, y)}{(u, v)} = \begin{vmatrix}
\pdv{x}{u} & \pdv{x}{v} \\[0.5em]
\pdv{y}{u} & \pdv{y}{v}
\end{vmatrix} = \pdv{x}{u}\pdv{y}{v} - \pdv{x}{v}\pdv{y}{u}
\]

For three dimensions $(x, y, z) \to (u, v, w)$:
\[
J = \begin{vmatrix}
\pdv{x}{u} & \pdv{x}{v} & \pdv{x}{w} \\[0.5em]
\pdv{y}{u} & \pdv{y}{v} & \pdv{y}{w} \\[0.5em]
\pdv{z}{u} & \pdv{z}{v} & \pdv{z}{w}
\end{vmatrix}
\]

\textbf{Example - Polar coordinates:} For $x = r\cos\theta$, $y = r\sin\theta$:
\[
J = \begin{vmatrix}
\pdv{x}{r} & \pdv{x}{\theta} \\[0.5em]
\pdv{y}{r} & \pdv{y}{\theta}
\end{vmatrix} = \begin{vmatrix}
\cos\theta & -r\sin\theta \\
\sin\theta & r\cos\theta
\end{vmatrix} = r\cos^2\theta + r\sin^2\theta = r
\]

Therefore: $dx \, dy = r \, dr \, d\theta$

This explains why integrals in polar coordinates have the extra factor of $r$:
\[
\iint_R f(x, y) \, dx \, dy = \iint_{R'} f(r\cos\theta, r\sin\theta) \, r \, dr \, d\theta
\]

\textbf{Differential forms:} A differential $df$ of a function $f(x, y)$ is:
\[
df = \pdv{f}{x}dx + \pdv{f}{y}dy
\]

This represents the infinitesimal change in $f$ due to infinitesimal changes $dx$ and $dy$ in the coordinates. More generally:
\[
df = \grad f \cdot d\vec{r}
\]
where $d\vec{r} = dx\hat{i} + dy\hat{j} + dz\hat{k}$ is the displacement vector.

\textbf{Einstein notation:} Using indexed coordinates $x^i$, the differential can be written compactly as:
\[
df = \pdv{f}{x^i}dx^i
\]
where summation over repeated index $i$ is implied. This notation emphasizes that the differential is a linear combination of coordinate differentials weighted by partial derivatives.

\subsection{Arc Length and Velocity Vector Tangent to a Curve}

A \textbf{parametric curve} in space is described by a position vector $\vec{r}(t)$ that depends on parameter $t$:
\[
\vec{r}(t) = x(t)\hat{i} + y(t)\hat{j} + z(t)\hat{k}
\]

\textbf{Velocity vector:} The derivative of position with respect to the parameter gives the velocity:
\[
\vec{v}(t) = \dv{\vec{r}}{t} = \dv{x}{t}\hat{i} + \dv{y}{t}\hat{j} + \dv{z}{t}\hat{k} = \dot{x}\hat{i} + \dot{y}\hat{j} + \dot{z}\hat{k}
\]

The velocity vector $\vec{v}(t)$ is \textbf{tangent to the curve} at each point and points in the direction of motion.

\textbf{Speed:} The magnitude of the velocity vector gives the speed:
\[
|\vec{v}(t)| = \sqrt{\left(\dv{x}{t}\right)^2 + \left(\dv{y}{t}\right)^2 + \left(\dv{z}{t}\right)^2}
\]

\textbf{Unit tangent vector:} Normalizing the velocity gives a unit vector tangent to the curve:
\[
\vec{T}(t) = \frac{\vec{v}(t)}{|\vec{v}(t)|}
\]

\textbf{Arc length:} The length of a curve from $t = a$ to $t = b$ is:
\[
s = \int_a^b |\vec{v}(t)| \, dt = \int_a^b \sqrt{\left(\dv{x}{t}\right)^2 + \left(\dv{y}{t}\right)^2 + \left(\dv{z}{t}\right)^2} \, dt
\]

In differential form:
\[
ds = |\vec{v}| \, dt = \sqrt{dx^2 + dy^2 + dz^2}
\]

\textbf{Arc length parameter:} We can reparametrize a curve using arc length $s$ as the parameter. Then:
\[
\dv{\vec{r}}{s} = \vec{T}
\]
is automatically a unit tangent vector (since $\left|\dv{\vec{r}}{s}\right| = 1$).

\textbf{Example:} Consider the helix $\vec{r}(t) = \cos(t)\hat{i} + \sin(t)\hat{j} + t\hat{k}$ for $0 \leq t \leq 2\pi$.

Velocity vector:
\[
\vec{v}(t) = -\sin(t)\hat{i} + \cos(t)\hat{j} + \hat{k}
\]

Speed:
\[
|\vec{v}(t)| = \sqrt{\sin^2(t) + \cos^2(t) + 1} = \sqrt{2}
\]

Arc length from $t = 0$ to $t = 2\pi$:
\[
s = \int_0^{2\pi} \sqrt{2} \, dt = 2\pi\sqrt{2}
\]

Unit tangent vector:
\[
\vec{T}(t) = \frac{1}{\sqrt{2}}\left(-\sin(t)\hat{i} + \cos(t)\hat{j} + \hat{k}\right)
\]

\textbf{Connection to tensor calculus:} In general curvilinear coordinates, the tangent vectors to coordinate curves form the basis vectors of the coordinate system. Understanding how tangent vectors transform is fundamental to tensor analysis.


\section{Basis vectors as partial derivatives and Jacobian}

In two-dimensional Euclidean space, a point is described by Cartesian coordinates
\[
(x,y)
\]

The position vector is
\[
\mathbf{R} = x \begingroup\color{blue}\vec{e_x}\endgroup + y \begingroup\color{blue}\vec{e_y}\endgroup
\]
where the basis vectors are
\[
\begingroup\color{blue}\vec{e_x}\endgroup = \begin{pmatrix}1 \\ 0\end{pmatrix},
\qquad
\begingroup\color{blue}\vec{e_y}\endgroup = \begin{pmatrix}0 \\ 1\end{pmatrix}.
\]

These satisfy:
\[
\begingroup\color{blue}\vec{e_x}\endgroup \cdot \begingroup\color{blue}\vec{e_x}\endgroup = 1, 
\quad
\begingroup\color{blue}\vec{e_y}\endgroup \cdot \begingroup\color{blue}\vec{e_y}\endgroup = 1,
\quad
\begingroup\color{blue}\vec{e_x}\endgroup \cdot \begingroup\color{blue}\vec{e_y}\endgroup = 0.
\]

Thus, $\{\begingroup\color{blue}\vec{e_x}\endgroup, \begingroup\color{blue}\vec{e_y}\endgroup\}$ is an orthonormal basis.

A key feature:

\[
\frac{\partial \begingroup\color{blue}\vec{e_x}\endgroup}{\partial x} = 
\frac{\partial \begingroup\color{blue}\vec{e_x}\endgroup}{\partial y} =
\frac{\partial \begingroup\color{blue}\vec{e_y}\endgroup}{\partial x} =
\frac{\partial \begingroup\color{blue}\vec{e_y}\endgroup}{\partial y} = 0.
\]

The basis vectors are constant everywhere in space.\\

We now define polar coordinates $(r,\theta)$ by the transformation

\[
x = r \cos\theta,
\qquad
y = r \sin\theta.
\]

The position vector becomes

\[
\mathbf{R} = r\cos\theta \,\begingroup\color{blue}\vec{e_x}\endgroup + r\sin\theta \,\begingroup\color{blue}\vec{e_y}\endgroup
\]

We define the radial unit vector as

\[
\begingroup\color{red}\tilde{\vec{e_r}}\endgroup
= \cos\theta \,\begingroup\color{blue}\vec{e_x}\endgroup 
+ \sin\theta \,\begingroup\color{blue}\vec{e_y}\endgroup.
\]

This points directly away from the origin.\\
And then define the angular unit vector as

\[
\begingroup\color{red}\tilde{\vec{e_\theta}}\endgroup
= -\sin\theta \,\begingroup\color{blue}\vec{e_x}\endgroup
+ \cos\theta \,\begingroup\color{blue}\vec{e_y}\endgroup
\]

This points in the direction of increasing $\theta$.\\
These vectors satisfy

\[
\begingroup\color{red}\tilde{\vec{e_r}}\endgroup \cdot \begingroup\color{red}\tilde{\vec{e_r}}\endgroup = 1,
\quad
\begingroup\color{red}\tilde{\vec{e_\theta}}\endgroup \cdot \begingroup\color{red}\tilde{\vec{e_\theta}}\endgroup = 1,
\quad
\begingroup\color{red}\tilde{\vec{e_r}}\endgroup \cdot \begingroup\color{red}\tilde{\vec{e_\theta}}\endgroup = 0
\]

So they form an orthonormal basis at each point, but polar basis vectors depend on $\theta$\\

Compute derivatives:

\[
\frac{\partial \begingroup\color{red}\tilde{\vec{e_r}}\endgroup}{\partial \theta}
=
-\sin\theta\,\begingroup\color{blue}\vec{e_x}\endgroup
+
\cos\theta\,\begingroup\color{blue}\vec{e_y}\endgroup
=
\begingroup\color{red}\tilde{\vec{e_\theta}}\endgroup
\]

\[
\frac{\partial \begingroup\color{red}\tilde{\vec{e_\theta}}\endgroup}{\partial \theta}
=
-\cos\theta\,\begingroup\color{blue}\vec{e_x}\endgroup
-
\sin\theta\,\begingroup\color{blue}\vec{e_y}\endgroup
=
-\begingroup\color{red}\tilde{\vec{e_r}}\endgroup
\]

Thus,

\[
\frac{\partial \begingroup\color{red}\tilde{\vec{e_r}}\endgroup}{\partial r} = 0,
\qquad
\frac{\partial \begingroup\color{red}\tilde{\vec{e_\theta}}\endgroup}{\partial r} = 0
\]

Key result:

\[
\frac{\partial \begingroup\color{red}\tilde{\vec{e_r}}\endgroup}{\partial \theta} = \begingroup\color{red}\tilde{\vec{e_\theta}}\endgroup,
\qquad
\frac{\partial \begingroup\color{red}\tilde{\vec{e_\theta}}\endgroup}{\partial \theta} = -\begingroup\color{red}\tilde{\vec{e_r}}\endgroup
\]

The basis vectors rotate as $\theta$ changes.\\

Now if you think about this, to find out the directions of the basis vectors, what we do is
following the directions of the coordinate lines and in the case of polar coordinates, for the $\begingroup\color{red}\tilde{\vec{e_\theta}}\endgroup$
basis vector, since the coordinate lines are actually curved, we need to pick the straightest possible
direction when moving along it, which results in being the vector tangent to the curve.\\

\textbf{This reasoning gives us a hint for the next important reinterpretation of basis vectors, as represented by a 
partial derivative of the position vector R}

\begin{equation}
    \begingroup\color{blue}\vec{e_x}\endgroup \equiv \frac{\partial \mathbf{R}}{\partial \begingroup\color{blue}x\endgroup}
\end{equation}

If you carefully think about it, it might look weird at first, but it will eventually make sense.
Think about $\mathbf{R}$ as a function of the coordinates x and y for example, and then
consider adding a small value $\Delta x$ to the x coordinate. The difference between these two position
vectors results in a vector that is parallel to the x coordinate axis.
So this resultant vector has the direction of the basis vector $\color{blue}\vec{e_x}$. \\
This observation is interesting, but we need to solve tha fact that the magnitude of this vector
is varying , depending on the value of the $\Delta x$ used. \\
This can be solved by dividing the resultant vector by the quantity $\Delta x$.
By doing so, we obtain a vector pointing in the x-basis vector direction and magnitude 1...now we 
can proceed further by taking the limit as $\Delta x$ approaches to 0:

\begin{equation}
    \frac{\partial \mathbf{R}}{\partial x} = \lim_{\Delta x \to 0}\frac{\mathbf{R(x + \Delta x, y)} - \mathbf{R(x,y)}}{\Delta x}
\end{equation}

Since this derivative always has a length of one and points the x-axis, we can regard it as the basis vector:

\begin{equation}
    \begingroup\color{blue}\vec{e_x}\endgroup \equiv \frac{\partial \mathbf{R}}{\partial \begingroup\color{blue}x\endgroup}
\end{equation}

This reasoning can be easily applied to the other direction, but generalized further
to any coordinate system.\\
The neat thing about this construction is that you can do it for any position vector, and you 
will still get your basis vectors. Fundamentally, this trick work because, in obtaining a slightly 
off position vector, a small change in coordinate is added in one coordinate only while keeping other coordinates constant. 
Then, we divide the resultant vector (obtained from the difference between the original position vector and the new slightly off one) 
by that small change to normalize the vector.\\\\

Now you should remember when we discussed, in the other notes for "Tensors for Beginners" about the 
forward and backward transformations when we move from one basis vector coordinate system to another.\\
Turns out that those transformation rules can be obtained easily once we introduced this above
concept about basis vectors being partial derivatives of the position vector with respect to the 
coordinate variables.\\
And also, that these transformation matrices have a proper name: \textbf{the Jacobian matrices}\\

Let's say we want to find the forward transformation matrix (the Jacobian) matric when moving 
from cartesian to polar coordinates. This would be just a matter of applying the chain rule:

\begin{align*}
    \begingroup\color{red}\tilde{\vec{e_r}}\endgroup &= \frac{\partial R}{\partial r} = \frac{\partial R}{\partial x}\frac{\partial x}{\partial r} + \frac{\partial R}{\partial y}\frac{\partial y}{\partial r} = \frac{\partial x}{\partial r}\begingroup\color{blue}\vec{e_x}\endgroup + \frac{\partial y}{\partial r}\begingroup\color{blue}\vec{e_y}\endgroup\\
    \begingroup\color{red}\tilde{\vec{e_\theta}}\endgroup &= \frac{\partial R}{\partial \theta} = \frac{\partial R}{\partial x}\frac{\partial x}{\partial \theta} + \frac{\partial R}{\partial y}\frac{\partial y}{\partial \theta} = \frac{\partial x}{\partial \theta}\begingroup\color{blue}\vec{e_x}\endgroup + \frac{\partial y}{\partial \theta}\begingroup\color{blue}\vec{e_y}\endgroup
\end{align*}

We know the transformation rules between cartesian and polar coordinates:

\[
x = r \cos\theta,
\qquad
y = r \sin\theta.
\]

so it's a matter of taking the derivatives to find that the Jacobian matrix, or forward matrix is:

\begin{align*}
    \frac{\partial R}{\partial r} &= \frac{\partial R}{\partial x}\frac{\partial x}{\partial r} + \frac{\partial R}{\partial y}\frac{\partial y}{\partial r} = cos\theta \frac{\partial R}{\partial x} + sin\theta \frac{\partial R}{\partial y}\\
    \frac{\partial R}{\partial \theta} &= \frac{\partial R}{\partial x}\frac{\partial x}{\partial \theta} + \frac{\partial R}{\partial y}\frac{\partial y}{\partial \theta} = -r\sin\theta \frac{\partial R}{\partial x} + r\cos\theta \frac{\partial R}{\partial y}
\end{align*}

So:

\begin{equation}
    F =
    \begin{bmatrix}
        \frac{\partial x}{\partial r} & \frac{\partial x}{\partial \theta}\\
        \frac{\partial y}{\partial r} & \frac{\partial y}{\partial \theta}
    \end{bmatrix} = 
    \begin{bmatrix}
        cos\theta & -r\sin\theta\\
        \sin\theta & r\cos\theta
    \end{bmatrix}
\end{equation}

And you can tackle, with the same derivation logic, the inverse Jacobian matrix, of backward transformation 
matrix, to obviously find out that the multiplication of the two yields the identity matrix.\\

Now, following a more tensorial approach, we can also re-write these transformation rules as:

\begin{align*}
    \frac{\partial \begingroup\color{violet}R\endgroup}{\partial \begingroup\color{magenta}p^i\endgroup} &= \frac{\partial \begingroup\color{violet}R\endgroup}{\partial \begingroup\color{cyan}c^j\endgroup}\frac{\partial \begingroup\color{cyan}c^j\endgroup}{\partial \begingroup\color{magenta}p^i\endgroup} \\
    \frac{\partial \begingroup\color{violet}R\endgroup}{\partial \begingroup\color{cyan}c^i\endgroup} &= \frac{\partial \begingroup\color{violet}R\endgroup}{\partial \begingroup\color{magenta}p^j\endgroup}\frac{\partial \begingroup\color{magenta}p^j\endgroup}{\partial \begingroup\color{cyan}c^i\endgroup} 
\end{align*}




\section{Vectors as derivatives and their transformation fules}

Allright we saw that basis vectors can be re-interpreted as partial derivative of the position
vector along the coordinate axis.\\
We can sort of extend this logic to the vectors themselves. We saw that figuring out vector 
components in a particular basis is kind of easy to do for individual vectors, but now we're going do it 
for vector fields (i.e. a vector for every point in space), \textbf{but for now, we're going to deal with a slightly 
different type of vector field, and that's a vector field along a curve, so basically we're defining 
a vector at every point along a curve.}\\

The way we imagine vector fields along a curve is considering the set of tangent vectors to the curves.
This gives us a vector field defined at every point on the curve.\\

A curve can be tought as a function that takes some input parameter $\lambda$ and outputs the position vector 
$\mathbf{R}$ that traces our curve in the space.\\
If we want to have a vector tangent to this curve, we just need to consider taking the derivative of the position 
vector with respect to its parameter:

\begin{equation}
    \frac{\dd \begingroup\color{violet}\vec{R}\endgroup}{\dd \begingroup\color{orange}\lambda\endgroup}
\end{equation}

Now you may think that finding out the components of this vector in a specific basis can be 
difficult but it is really not, if we make use of the multivariable chain rule and expand it out in the
coodinates we're working in.\\
For example, in the cartesian coordinates $\color{cyan}c^i$ and using the Einstein summation convention notation:

\begin{equation}
    \frac{\dd \begingroup\color{violet}\vec{R}\endgroup}{\dd \begingroup\color{orange}\lambda\endgroup} = 
    \frac{\dd \begingroup\color{cyan}c^i\endgroup}{\dd \begingroup\color{orange}\lambda\endgroup}
    \frac{\partial \begingroup\color{violet}\vec{R}\endgroup}{\partial \begingroup\color{cyan}c^i\endgroup}
\end{equation}

At the end of the day, considering doing it in cartesian $\mathbb{R^2}$ for simplicity, these are equivalent (individual vector on the top, and vector field along a curve on the bottom):

\begin{align*}
    \begingroup\color{orange}\vec{v}\endgroup &= \begingroup\color{cyan}v^1\endgroup \begingroup\color{blue}\vec{e_1}\endgroup + \begingroup\color{cyan}v^2\endgroup \begingroup\color{blue}\vec{e_2}\endgroup\\
    \frac{\dd \begingroup\color{violet}\vec{R}\endgroup}{\dd \begingroup\color{orange}\lambda\endgroup} &= \frac{\dd \begingroup\color{cyan}x\endgroup}{\dd \begingroup\color{orange}\lambda\endgroup}
    \frac{\partial \begingroup\color{violet}\vec{R}\endgroup}{\partial \begingroup\color{cyan}x\endgroup} + 
    \frac{\dd \begingroup\color{cyan}y\endgroup}{\dd \begingroup\color{orange}\lambda\endgroup}
    \frac{\partial \begingroup\color{violet}\vec{R}\endgroup}{\partial \begingroup\color{cyan}y\endgroup}
\end{align*}

The same can be applied to whatever coordinate system we're working in, so to give you an example, let's say we
have the circular curve of radius 2, centered at the origin in $\mathbb{R^2}$, defined by:

\begin{equation}
    \begingroup\color{violet}\vec{R}\endgroup (\begingroup\color{orange}\lambda\endgroup) = \left(\begingroup\color{cyan}x\endgroup(\begingroup\color{orange}\lambda\endgroup), \begingroup\color{cyan}y\endgroup(\begingroup\color{orange}\lambda\endgroup)\right) =
    \left(2 \cos\lambda, 2 \sin\lambda\right)
\end{equation}

The vector field along the curve is made up of the tangent vectors along the curve, that 
can be expanded as defined:

\begin{align*}
    \frac{\dd \begingroup\color{violet}\vec{R}\endgroup}{\dd \begingroup\color{orange}\lambda\endgroup} &= \frac{\dd \begingroup\color{cyan}x\endgroup}{\dd \begingroup\color{orange}\lambda\endgroup}
    \frac{\partial \begingroup\color{violet}\vec{R}\endgroup}{\partial \begingroup\color{cyan}x\endgroup} + 
    \frac{\dd \begingroup\color{cyan}y\endgroup}{\dd \begingroup\color{orange}\lambda\endgroup}
    \frac{\partial \begingroup\color{violet}\vec{R}\endgroup}{\partial \begingroup\color{cyan}y\endgroup}\\
    &= -2 \sin\lambda \frac{\partial \begingroup\color{violet}\vec{R}\endgroup}{\partial \begingroup\color{cyan}x\endgroup} + 
    2 \cos\lambda\frac{\partial \begingroup\color{violet}\vec{R}\endgroup}{\partial \begingroup\color{cyan}y\endgroup}
\end{align*}

The components can also be written, for the cartesian basis, as a vector column:

\begin{equation}
    \begin{bmatrix}
        -2\sin\lambda\\
        2 \cos\lambda
    \end{bmatrix}_{\frac{\partial}{\partial \begingroup\color{cyan}c^i\endgroup}}
\end{equation}

We can do the same in polar coordinates $r, \theta$, where the same curve is defined as:

\begin{equation}
    \begingroup\color{violet}\vec{R}\endgroup (\begingroup\color{orange}\lambda\endgroup)= \left(\begingroup\color{magenta}r\endgroup(\begingroup\color{orange}\lambda\endgroup), \begingroup\color{magenta}\theta\endgroup(\begingroup\color{orange}\lambda\endgroup)\right) =
    \left(2 , \lambda\right)
\end{equation}

\begin{align*}
    \frac{\dd \begingroup\color{violet}\vec{R}\endgroup}{\dd \begingroup\color{orange}\lambda\endgroup} &= \frac{\dd \begingroup\color{magenta}r\endgroup}{\dd \begingroup\color{orange}\lambda\endgroup}
    \frac{\partial \begingroup\color{violet}\vec{R}\endgroup}{\partial \begingroup\color{magenta}r\endgroup} + 
    \frac{\dd \begingroup\color{magenta}\theta\endgroup}{\dd \begingroup\color{orange}\lambda\endgroup}
    \frac{\partial \begingroup\color{violet}\vec{R}\endgroup}{\partial \begingroup\color{magenta}\theta\endgroup}\\
    &= 0 \frac{\partial \begingroup\color{violet}\vec{R}\endgroup}{\partial \begingroup\color{magenta}r\endgroup} + 
    1 \frac{\partial \begingroup\color{violet}\vec{R}\endgroup}{\partial \begingroup\color{magenta}\theta\endgroup}
\end{align*}

So:

\begin{equation}
    \begin{bmatrix}
        0\\
        1
    \end{bmatrix}_{\frac{\partial}{\partial \begingroup\color{magenta}p^i\endgroup}}
\end{equation}

Now, as you can already imagine, like we saw that for individual vectors, the components are contravariant (i.e. they transform the opposite way of the basis vectors), 
we're going to simply show that this is valid indeed in this vector fields case.\\

We already saw in the previous section, the transformation rules of the basis vectors using the Jacobain and Inverse Jacobian:

\begin{subequations}\label{eq:forward-backward}
\begin{empheq}[box=\widefbox]{align}
    \frac{\partial \begingroup\color{violet}\vec{R}\endgroup}{\partial \begingroup\color{magenta}p^i\endgroup} &= \frac{\partial \begingroup\color{violet}\vec{R}\endgroup}{\partial \begingroup\color{cyan}c^j\endgroup}\frac{\partial \begingroup\color{cyan}c^j\endgroup}{\partial \begingroup\color{magenta}p^i\endgroup} & \text{Forward Transform}\\
    \frac{\partial \begingroup\color{violet}\vec{R}\endgroup}{\partial \begingroup\color{cyan}c^i\endgroup} &= \frac{\partial \begingroup\color{violet}\vec{R}\endgroup}{\partial \begingroup\color{magenta}p^j\endgroup}\frac{\partial \begingroup\color{magenta}p^j\endgroup}{\partial \begingroup\color{cyan}c^i\endgroup} & \text{Backward Transform}
\end{empheq}
\end{subequations}

Now we can just use these and replace the basis vectors with their new basis transformation rule: 

\begin{equation}
    \frac{\dd \begingroup\color{violet}\vec{R}\endgroup}{\dd \begingroup\color{orange}\lambda\endgroup} = 
    \frac{\dd \begingroup\color{cyan}c^i\endgroup}{\dd \begingroup\color{orange}\lambda\endgroup}
    \frac{\partial \begingroup\color{violet}\vec{R}\endgroup}{\partial \begingroup\color{cyan}c^i\endgroup} = 
    \frac{\dd \begingroup\color{cyan}c^i\endgroup}{\dd \begingroup\color{orange}\lambda\endgroup}
    \left(\frac{\partial \begingroup\color{violet}\vec{R}\endgroup}{\partial \begingroup\color{magenta}p^j\endgroup}\frac{\partial \begingroup\color{magenta}p^j\endgroup}{\partial \begingroup\color{cyan}c^i\endgroup}\right) = 
    \left(\frac{\dd \begingroup\color{cyan}c^i\endgroup}{\dd \begingroup\color{orange}\lambda\endgroup} \frac{\partial \begingroup\color{magenta}p^j\endgroup}{\partial \begingroup\color{cyan}c^i\endgroup}\right)
    \frac{\partial \begingroup\color{violet}\vec{R}\endgroup}{\partial \begingroup\color{magenta}p^j\endgroup}
\end{equation}

This is a linear combination of the new basis vector $\color{magenta}p^j$, and since we know that we can write:

\begin{equation}
    \frac{\dd \begingroup\color{violet}\vec{R}\endgroup}{\dd \begingroup\color{orange}\lambda\endgroup} = 
    \frac{\dd \begingroup\color{magenta}p^i\endgroup}{\dd \begingroup\color{orange}\lambda\endgroup}
    \frac{\partial \begingroup\color{violet}\vec{R}\endgroup}{\partial \begingroup\color{magenta}p^i\endgroup}
\end{equation}

This means that:

\begin{equation}
    \frac{\dd \begingroup\color{magenta}p^i\endgroup}{\dd \begingroup\color{orange}\lambda\endgroup} = 
    \frac{\dd \begingroup\color{cyan}c^i\endgroup}{\dd \begingroup\color{orange}\lambda\endgroup} \frac{\partial \begingroup\color{magenta}p^j\endgroup}{\partial \begingroup\color{cyan}c^i\endgroup}
\end{equation}

Which shows that the vector field components transform using the backward transformation (inverse Jacobian)\\

We can actually test this out in the previous example right?
So we have the components being, in both coordinate systems:

\begin{align*}
    \begin{bmatrix}
        -2\sin\lambda\\
        2 \cos\lambda
    \end{bmatrix}_{\frac{\partial }{\partial \begingroup\color{cyan}c^i\endgroup}} &&
    \begin{bmatrix}
        0\\
        1
    \end{bmatrix}_{\frac{\partial }{\partial \begingroup\color{magenta}p^i\endgroup}}
\end{align*}

And we also know the Inverse Jacobian in the case of polar-to-cartesian coordinates, being:

\begin{equation}
    B = J^{-1} =
    \begin{bmatrix}
        x/r & y/r\\
        -y/r^2 & x/r^2
    \end{bmatrix} = 
    \begin{bmatrix}
        \cos\lambda & \sin\lambda\\
        -1/2\sin\lambda & 1/2\cos\lambda
    \end{bmatrix}
\end{equation}

Let's check that:

\begin{align*}
    B \begin{bmatrix}
        -2\sin\lambda\\
        2 \cos\lambda
    \end{bmatrix}_{\frac{\partial }{\partial \begingroup\color{cyan}c^i\endgroup}} &= 
     \begin{bmatrix}
        \cos\lambda & \sin\lambda\\
        -1/2\sin\lambda & 1/2\cos\lambda
    \end{bmatrix} 
    \begin{bmatrix}
        -2\sin\lambda\\
        2 \cos\lambda
    \end{bmatrix}_{\frac{\partial }{\partial \begingroup\color{cyan}c^i\endgroup}}\\ 
    &= \begin{bmatrix}
        -\cos\lambda 2\sin\lambda + \sin\lambda 2\cos\lambda\\
        -1/2 \sin\lambda (-2 \sin\lambda) + 1/2 \cos\lambda 2 \cos\lambda
    \end{bmatrix}_{\frac{\partial}{\partial \begingroup\color{magenta}p^i\endgroup}}\\
    & = \begin{bmatrix}
        0\\
        1
    \end{bmatrix}_{\frac{\partial}{\partial \begingroup\color{magenta}p^i\endgroup}}
\end{align*}

Same can be checked on the opposite direction applying the forward/jacobian transformation.\\

One last thing to mention before terminating this section is in terms of notation: you notice that 
the position vector appears in a lot of them, so what we're going to do from here onward, is get rid of the 
position vector symbol.\\
This might seem like a weird choice of notation, but from now on, for example, the $\mathbb{R}^2$ cartesian basis vectors will be 
written as:

\begin{align*}
    \begingroup\color{blue}\vec{e_x}\endgroup = \frac{\partial}{\partial \begingroup\color{cyan}x\endgroup}\\
    \begingroup\color{blue}\vec{e_y}\endgroup = \frac{\partial}{\partial \begingroup\color{cyan}y\endgroup}
\end{align*}

\textbf{where basically we consider basis vectors these partial derivative operators.}\\
If you think about this, it's not very weird at the end, as derivative operators are linear, as they can be scaled and added like vectors do, 
so they have both a magnitude and a direction.
They also obey the forward and backward transformation rules that we would expect vectors to obey.\\

The main logical reason to get rid of the position vector symbol is that, the position vector definition requires an origin to be set, and 
later on when we're going to deal with curved surfaces and manifolds, it's not easy to define an origin point.\\
So, using a position vector without an origin does not make a lot of sense, but we're still going to be able to use
vectors and vector fields as long as we have these derivative operators to work with.\\






% -------------------- bibliography --------------------
\begin{thebibliography}{9}
\bibitem{eigenchris}
eigenchris, \textit{Tensors Calculus}, YouTube video series,\\
\url{https://www.youtube.com/playlist?list=PLJHszsWbB6hpk5h8lSfBkVrpjsqvUGTCx}
\end{thebibliography}

\end{document}
